\begin{statementbox}
	\textbf{\textcolor{CStatement}{条件}}
	\begin{description}[style=nextline, leftmargin=2.8em, labelsep=0.5em, itemsep=0.35em]
		\item[\textbf{\textcolor{CStatement}{条件1（点在线上）：}}] 点$P(x_0,y_0)$在抛物线$C: y^2=2px$（$p>0$）上，即$y_0^2=2px_0$成立。
	\end{description}

	\textbf{\textcolor{CStatement}{结论}}
	\begin{description}[style=nextline, leftmargin=2.8em, labelsep=0.5em, itemsep=0.45em]
		\item[\textbf{\textcolor{CStatement}{结论一（切线方程）：}}]
			\textbf{\textcolor{CStatement}{直观描述：}}
			切线方程可视为将原方程中的$y^2$换为$y_0y$，$2px$换为$p(x+x_0)$。
			\[
				y_0 y = p(x + x_0).
			\]

		\item[\textbf{\textcolor{CStatement}{推广结论（其他开口）：}}]
			\textbf{\textcolor{CStatement}{直观描述：}}
			对其他方向的标准抛物线，切线公式可由类似替换得到。

			（1）开口向左（$y^2=-2px$）：
			\[
				y_0 y = -p(x + x_0).
			\]
			（2）开口向上（$x^2=2py$）：
			\[
				x_0 x = p(y + y_0).
			\]
			（3）开口向下（$x^2=-2py$）：
			\[
				x_0 x = -p(y + y_0).
			\]
	\end{description}
\end{statementbox}
